\section{REFERENCIAL TEÓRICO}
Destina-se a apresentar as leituras e os fundamentos teóricos que embasam a proposta da pesquisa. O pesquisador deve se atentar para o fato de que citação, conforme \textcite{abnt10520}, de até três linhas não recebe recuo, às maiores que isso deve ser dado um recuo de 4 cm da margem esquerda do texto.

Usar fonte Arial ou Times New Roman, tamanho 10 ou 11, nestes casos: nas citações com mais de três linhas, nas notas de rodapé, na indicação da fonte (autoria) das figuras, quadros e tabelas, bem como nas legendas e nos conteúdos dos quadros e das tabelas. Para a citação com mais de três linhas deixar uma linha em branco antes e outra após a citação direta e recuo de 4 cm da margem esquerda. (Autor, ano, p. xx).

Vale ressaltar que caso a citação seja oriunda de um trabalho em língua estrangeira, cuja tradução foi feita pelo autor do TCC, essa informação deve vir também dentro dos parênteses por meio da expressão “tradução nossa”, conforme exemplificado a seguir:

\begin{quote}
    “citação traduzida citação traduzida citação traduzida citação traduzida citação traduzida.” (Autor, ano, p. xx, tradução nossa).
\end{quote}

Nos casos em que a fonte consultada da citação não possui paginação, a nota de rodapé deve trazer a seguinte informação: A fonte consultada não é paginada. Nesse caso, se enquadram informações extraídas de sites do governo, por exemplo, conforme exemplificado na NBR 10520/2023, item 7.1, exemplo 6:

\begin{quote}
    Exemplo 6: “O Poder Executivo envidará esforços no sentido de antecipar a entrega do Plano previsto no caput deste artigo em pelo menos 15 dias (Brasil, 1999).
    
    Nota: A fonte consultada neste exemplo 6 não é paginada.”
\end{quote}

\subsection{TÍTULO DA SEÇÃO SECUNDÁRIA}
O autor deve nomear as seções de acordo com a sua pesquisa.

\subsubsection{Título da seção terciária}
Texto texto texto texto texto texto texto texto texto texto texto texto texto texto texto texto texto texto texto texto texto texto texto texto texto texto texto texto texto texto texto texto texto texto texto texto texto texto texto texto texto.

\paragraph{Título da seção quaternária}
Texto texto texto texto texto texto texto texto texto texto texto texto texto texto texto texto texto texto texto texto texto.